\section{Spectral family of a bounded self-adjoint linear operator}

2
3
4
The spectral family is pretty abstract, the following exercises are meant to compute it in a concrete case:
5
6
7 Also compute the $E_lambda$ !
Extra: Take any function $f:[0,1]->R$, let's say $f(x) = x - 1/2$, and compute the same quantities as in 5,6,7 and the $E_lambda$, for the operator $T:L^2 \to L^2: g \to fg$
8 Explains the name $T^+$
9
10 Two more computations in trivial cases

\begin{exercise}{2}
Suppose $E_\lambda$ satisfies all conditions to be a spectral family (9.7-1) except for $E_{\lambda + 0} x = \lim_{\mu \to \lambda + 0} E_\mu x = E_\lambda x$.
Find $\tilde{E}_\lambda$ satisfying all those conditions, including the limit condition.
\end{exercise}
\begin{proof}
fill
\end{proof}

\begin{exercise}{3}
Prove that $T^- T = T T^-$.
\end{exercise}
\begin{proof}
fill
\end{proof} 

\begin{exercise}{4}
Find $T^+, T^-, (T^2)^{1/2}$ and the other square roots of $T^2$ if
\begin{align*}
    T = \begin{bmatrix}
        2 & 0  \\
        0 & -3
    \end{bmatrix},
\end{align*}
\end{exercise}
\begin{proof}
fill
\end{proof} 

\begin{exercise}{5}
If in the finite dimensiona lcase a linear operator $T$ is represented by a real diagonal matrix $\tilde{T}$, what is the spectrum of $T$?
How do we obtain $\tilde{T}$ from the matrix
a) $\tilde{T}^+$ (representing $T^+$,
b) $\tilde{T}^-$ (representing $T^-$,
c) $\tilde{B}$ (representing $B$)?
\end{exercise}
\begin{proof}
fill
\end{proof} 

\begin{exercise}{6}
In exercise 5, how can we obtain the matrix representing the projection of $H$
a) onto $\cN(T^+)$,
b) onto $\cN(T^+_\lambda)$?
\end{exercise}
\begin{proof}
fill
\end{proof} 

\begin{exercise}{7}
Obtain from $\tilde{T}$ in exercise 5 the matrices representing
a) $T_\lambda$,
b) $T^+_\lambda$,
c) $T^-_\lambda$,
d) $B_\lambda$, 
e) $E_\lambda$.
\end{exercise}
\begin{proof}
fill
\end{proof} 

\begin{exercise}{Extra}
Take any function $f: [0,1] \to \R$, let's say $f(x) = x - 1/2$, and compute the same quantities as in 5, 6, 7 and the $E_\lambda$, for the operator $T:L² \to L²: g \to fg$.
\end{exercise}
\begin{proof}
fill
\end{proof} 

\begin{exercise}{8}
Show that if a bounded self-adjoint linear operator $T: H \to H$ is positive, then $T = T^+$ and $T^- = 0$.
\end{exercise}
\begin{proof}
fill
\end{proof} 

\begin{exercise}{9}
Find the spectral family of the zero operator $T = 0: H \to H$, where $H \neq \set{0}$.
\end{exercise}
\begin{proof}
fill
\end{proof} 

\begin{exercise}{10}
Let $T = I: H \to H$.
Find $B_\lambda = (T^2_\lambda)^{1/2}$, $T^+_\lambda$, $\cN(T^+_\lambda)$ and $E_\lambda$.
\end{exercise}
\begin{proof}
fill
\end{proof} 
